\documentclass[../../main.tex]{subfiles}

\begin{document}

\section{Hahn-Banach 定理}\label{section:ch14-Hahn-Banach}

\begin{definition}[正齐次、次可加]\label{definition:ch14-正齐次次可加}
线性空间 $X$ 上的泛函 $p:X\to[0,\infty)$ 称为是正齐次的, 若
\begin{align*}
\text{对所有 }x\in X,\ \lambda>0,\quad p(\lambda x)=\lambda p(x).
\end{align*}
称为是次可加的, 若
\begin{align*}
\text{对所有 }x,\ y\in X,\quad p(x+y)\leqslant p(x)+p(y).
\end{align*}
\end{definition}

线性空间上的任何范数既是次可加的(三角不等式)又是正齐次的.

\begin{lemma}[Hahn-Banach 引理]\label{lemma:ch14-HB引理}
令 $p$ 为线性空间 $X$ 上的正齐次、次可加的泛函而 $Y$ 是 $X$ 的子空间, 存在定义在其上的线性泛函 $\psi$, 满足
\begin{align*}
\text{在 }Y\text{ 上},\quad\psi\leqslant p.
\end{align*}
令 $z$ 属于 $X\sim Y$. 则 $\psi$ 可被延拓为 $\operatorname{span}[Y+z]$ 上的线性泛函 $\psi$, 使得
\begin{align*}
\text{在 }\operatorname{span}[Y+z]\text{ 上},\quad\psi\leqslant p.
\end{align*}
\end{lemma}

\begin{proof}
由于 $\operatorname{span}[Y+z]$ 中的每个向量可唯一地表示为 $y+\lambda z$, 其中 $y\in Y$ 而 $\lambda\in\mathbf{R}$, 仅须找到数 $\psi(z)$ 具有以下性质:
\begin{align}
\text{对所有 }y\in Y\text{ 与 }\lambda\in\mathbf{R},\quad\psi(y)+\lambda\psi(z)\leqslant p(y+\lambda z).\label{align:ch14-HB延拓}
\end{align}
的确, 对这样的数 $\psi(z)$, 对所有 $Y$ 中的 $y$ 与 $\lambda\in\mathbf{R}$ 定义 $\psi(y+\lambda z)=\psi(y)+\lambda\psi(z)$ 就得到所要的延拓.

对任何向量 $y_{1},y_{2}\in Y$, 由于 $\psi$ 是线性的, 在 $Y$ 上 $\psi\leqslant p$ 且 $p$ 是次可加的,
\begin{align*}
\psi(y_{1})+\psi(y_{2})=\psi(y_{1}+y_{2})\leqslant p(y_{1}+y_{2})\\
=p((y_{1}-z)+(y_{2}+z))\leqslant p(y_{1}-z)+p(y_{2}+z).
\end{align*}
因此,
\begin{align*}
\psi(y_{1})-p(y_{1}-z)\leqslant-\psi(y_{2})+p(y_{2}+z).
\end{align*}
当 $y_{1}$ 与 $y_{2}$ 在 $Y$ 中的所有向量变动, 该不等式左边的任何数不大于右边的任何数. 根据 $\mathbf{R}$ 的完备性, 若我们定义 $\psi(z)$ 为该不等式左边的数的上确界, 则 $\psi(z)\in\mathbf{R}$. 此外, 对任何 $y\in Y$, 根据 $\psi(z)$ 作为上界的选取, $\psi(y)-p(y-z)\leqslant\psi(z)$, 由 $\psi(z)$ 作为最小上界的选取, $\psi(z)\leqslant-\psi(y)+p(y+z)$. 因此对所有 $y\in Y$,
\begin{align}
\psi(y)-p(y-z)\leqslant\psi(z)\leqslant-\psi(y)+p(y+z).\label{align:ch14-HB界}
\end{align}
令 $y$ 属于 $Y$. 对 $\lambda>0$, 在不等式 $\psi(z)\leqslant-\psi(y)+p(y+z)$ 中用 $y/\lambda$ 代替 $y$, 用 $\lambda$ 乘以每边, 且用 $p$ 与 $\psi$ 的正齐次性得到想要的不等式\eqref{align:ch14-HB延拓}. 对 $\lambda<0$, 在不等式 $\psi(y)-p(y-z)\leqslant\psi(z)$ 中用 $-y/\lambda$ 代替 $y$, 用 $-\lambda$ 乘以每边, 再次用正齐次性得到想要的不等式\eqref{align:ch14-HB延拓}. 因此, 若选取数 $\psi(z)$ 使得\eqref{align:ch14-HB界}成立, 则\eqref{align:ch14-HB延拓}成立.
\end{proof}

\begin{theorem}[Hahn-Banach 定理]\label{theorem:ch14-HB}
令 $p$ 为线性空间 $X$ 上的正齐次、次可加的泛函而 $Y$ 是 $X$ 的子空间, 使得在其上有线性泛函 $\psi$ 满足
\begin{align*}
\text{在 }Y\text{ 上},\quad\psi\leqslant p.
\end{align*}
则 $\psi$ 可延拓为整个 $X$ 上的线性泛函 $\psi$, 使得在整个 $X$ 上 $\psi\leqslant p$.
\end{theorem}

\begin{proof}
考虑定义在 $X$ 的子空间 $Y_{\eta}$ 上的所有线性泛函 $\eta$ 所组成的族 $\mathcal{F}$, 使得 $Y\subseteq Y_{\eta}$, 在 $Y$ 上 $\eta=\psi$, 在 $Y_{\eta}$ 上 $\eta\leqslant p$. 由 $\psi$ 的延拓组成的这个特定的族 $\mathcal{F}$ 可通过以下方式定义偏序而成为偏序集: 若 $Y_{\eta_{1}}\subseteq Y_{\eta_{2}}$, $\eta_{1}\prec\eta_{2}$, 且在 $Y_{\eta_{1}}$ 上 $\eta_{1}=\eta_{2}$.

令 $\mathcal{F}_{0}$ 为 $\mathcal{F}$ 的全序子族. 定义 $Z$ 为 $\mathcal{F}_{0}$ 中的泛函的定义域的并. 由于 $\mathcal{F}_{0}$ 是全有序的, 任何两个这样的定义域中一个包含于另一个, 而由于每个定义域是 $X$ 的线性子空间, $Z$ 也是. 对 $z\in Y$, 选取 $\eta\in\mathcal{F}_{0}$ 使得 $z\in Y_{\eta}$: 定义 $\eta^{*}(z)=\eta(z)$. 则根据 $\mathcal{F}_{0}$ 的全序性, $\eta^{*}$ 是 $Z$ 上恰当定义的线性泛函. 观察到在 $Z$ 上 $\eta^{*}\leqslant p$, $Y\subseteq Z$ 且在 $Y$ 上 $\eta^{*}=\psi$, 由于 $\mathcal{F}_{0}$ 中的每个泛函具有这三个性质. 因此对所有 $\eta\in\mathcal{F}$, $\eta\prec\eta^{*}$. 因此 $\mathcal{F}$ 的每个全序子族有上界. 因此, 根据 Zorn 引理, $\mathcal{F}$ 有最大元 $\psi_{0}$. 令 $\psi_{0}$ 的定义域为 $Y_{0}$. 由定义, $Y\subseteq Y_{0}$ 且在 $Y_{0}$ 上 $\psi_{0}\leqslant p$. 我们从 Hahn-Banach 引理推出最大延拓 $\psi_{0}$ 定义在整个 $X$ 上.
\end{proof}

\begin{theorem}\label{theorem:ch14-范数延拓}
令 $X_{0}$ 为赋范线性空间 $X$ 的线性子空间. 则 $X_{0}$ 上的每个有界线性泛函 $\psi$ 可延拓为整个 $X$ 上的有界线性泛函, 延拓后的泛函与 $\psi$ 有相同的范数. 特别地, 对每个 $x\in X$, 存在 $\psi\in X^{*}$ 使得
\begin{align}
\psi(x)=\|x\|\text{ 且 }\|\psi\|=1.\label{align:ch14-范数延拓}
\end{align}
\end{theorem}

\begin{proof}
令 $\psi:X_{0}\to\mathbf{R}$ 为线性与有界的. 定义
\begin{align*}
M=\|\psi\|=\sup\{|\psi(x)|\mid x\in X_{0},\|x\|\leqslant 1\}.
\end{align*}
定义 $p:X\to\mathbf{R}$ 为
\begin{align*}
\text{对所有 }x\in X,\quad p(x)=M\|x\|.
\end{align*}
泛函 $p$ 是次可加的与正齐次的. 根据 $M$ 的定义,
\begin{align*}
\text{在 }X_{0}\text{ 上},\quad\psi\leqslant p.
\end{align*}
根据 Hahn-Banach 定理, $\psi$ 可以延拓为整个 $X$ 上的连续线性泛函 $\psi$, 且对所有 $x\in X$, $\psi(x)\leqslant p(x)=M\|x\|$. 用 $-x$ 代替 $x$, 我们推出对所有 $x\in X$, $|\psi(x)|\leqslant p(x)=M\|x\|$, 因此 $\psi$ 在整个 $X$ 上的延拓有与 $\psi:X_{0}\to\mathbf{R}$ 相同的范数.

现在令 $x\neq 0$ 属于 $X$. 定义 $\eta:\operatorname{span}[x]\to\mathbf{R}$ 为 $\eta(\lambda x)=\lambda\|x\|$. 观察到 $\|\eta\|=1$. 注意到 $\eta(x)\neq 0$. 根据证明的第一部分, 泛函 $\eta$ 可延拓为整个 $X$ 上的有界线性泛函, 延拓后的泛函也具有范数 $1$.
\end{proof}

\begin{example}\label{example:ch14-延拓C[a,b]}
令 $x_{0}$ 属于闭有界区间 $[a,b]$. 对所有 $f\in C[a,b]$ 定义
\begin{align*}
\psi(f)=f(x_{0}).
\end{align*}
我们考虑 $C[a,b]$ 为 $L^{\infty}[a,b]$ 的子空间(见习题 29). 我们从前一个定理推出 $\psi$ 可延拓为有界线性泛函 $\psi:L^{\infty}[a,b]\to\mathbf{R}$. 这样的延拓从未直接给出.
\end{example}

\begin{example}\label{example:ch14-Banach极限}
定义 $\ell^{\infty}$ 上的正齐次、次可加泛函 $p$ 为
\begin{align*}
\text{对所有 }\{x_{n}\}\in\ell^{\infty},\quad p(\{x_{n}\})=\limsup\{x_{n}\}.
\end{align*}
令 $c\subseteq\ell^{\infty}$ 为收敛序列的子空间. 定义 $c$ 上的 $L$ 为
\begin{align*}
\text{对所有 }\{x_{n}\}\in c,\quad L(\{x_{n}\})=\lim_{n\to\infty}x_{n}.
\end{align*}
由于 $L$ 是线性的且在 $c$ 上 $L\leqslant p$. $L$ 可延拓为 $\ell^{\infty}$ 上的线性泛函, 使得在 $\ell^{\infty}$ 上 $L\leqslant p$. 任何这样的延拓称为 Banach 极限.
\end{example}

在前一章我们考虑 Banach 空间的闭子空间 $X_{0}$ 是否有闭线性补. 以下推论告诉我们若 $X_{0}$ 是有限维的, 它确实有.

\begin{corollary}\label{corollary:ch14-有限维闭补}
令 $X$ 为赋范线性空间. 若 $X_{0}$ 是 $X$ 的有限维子空间, 则存在 $X$ 的闭线性子空间 $X_{1}$, 使得 $X=X_{0}\oplus X_{1}$.
\end{corollary}

\begin{proof}
令 $e_{1},\cdots,e_{n}$ 为 $X_{0}$ 的基. 对 $1\leqslant k\leqslant n$, 定义 $\psi_{k}:X_{0}\to\mathbf{R}$ 为 $\psi_{k}\left(\sum_{i=1}^{n}\lambda_{i}\cdot e_{i}\right)=\lambda_{k}$. 由于 $X_{0}$ 是有限维的, 这些 $\psi_{k}$ 是连续的. 根据定理 7, 每个 $\psi_{k}$ 有到整个 $X$ 上的延拓 $\psi_{k}'$. 因此每个 $\psi_{k}'$ 有闭的核, 使得子空间 $X_{1}=\bigcap_{k=1}^{n}\ker\psi_{k}'$ 也是闭的. 容易检验 $X=X_{0}\oplus X_{1}$.
\end{proof}

\begin{corollary}\label{corollary:ch14-J等距}
令 $X$ 为赋范线性空间. 则自然嵌入 $J:X\to X^{**}$ 是一个等距.
\end{corollary}

\begin{proof}
令 $x$ 属于 $X$. 观察到根据对偶空间上范数的定义,
\begin{align*}
\text{对所有 }\psi\in X^{*},\quad|\psi(x)|\leqslant\|\psi\|\cdot\|x\|.
\end{align*}
因此
\begin{align*}
\text{对所有 }\psi\in X^{*},\quad|J(x)(\psi)|\leqslant\|x\|\cdot\|\psi\|.
\end{align*}
因此 $J(x)$ 是有界的且 $\|J(x)\|\leqslant\|x\|$. 另一方面, 根据定理 7, 存在 $\psi\in X^{*}$ 使得 $\|\psi\|=1$ 且 $J(x)(\psi)=\|x\|$. 因此 $\|x\|\leqslant\|J(x)\|$. 我们得出 $J$ 是一个等距.
\end{proof}

\begin{theorem}\label{theorem:ch14-闭包判据}
令 $X_{0}$ 为赋范线性空间 $X$ 的子空间. 那么 $X$ 中的点 $x$ 属于 $X_{0}$ 的闭包当且仅当只要泛函 $\psi\in X^{*}$ 在 $X_{0}$ 上消失则它也在 $x$ 消失.
\end{theorem}

\begin{proof}
显然由连续性可知, 若连续泛函在 $X_{0}$ 上消失, 则它在 $X_{0}$ 的闭包上也消失. 为证明反命题, 令 $x_{0}$ 属于 $X\sim\overline{X_{0}}$. 我们必须证明存在 $\psi\in X^{*}$ 在 $X_{0}$ 上消失但 $\psi(x_{0})\neq 0$. 定义 $Z=\overline{X_{0}}\oplus[x_{0}]$ 与 $\psi:Z\to\mathbf{R}$ 为:
\begin{align*}
\text{对所有 }x\in\overline{X_{0}}\text{ 与 }\lambda\in\mathbf{R},\quad\psi(x+\lambda x_{0})=\lambda.
\end{align*}
我们宣称 $\psi$ 是有界的. 事实上, 由于 $\overline{X_{0}}$ 是闭的, 它的补是开的. 因此存在 $r>0$ 使得对所有 $u\in\overline{X_{0}}$, $\|u-x_{0}\|\geqslant r$. 因此, 对 $x\in\overline{X_{0}}$ 与 $\lambda\in\mathbf{R}$,
\begin{align*}
\|x+\lambda x_{0}\|=\|\lambda\|\cdot\|(-1/\lambda\cdot x)-x_{0}\|\geqslant|\lambda|\cdot r.
\end{align*}
从这我们推出 $\psi:X\to\mathbf{R}$ 是有界的且 $\|\psi\|\leqslant 1/r$. 定理 7 告诉我们 $\psi$ 可延拓为整个 $X$ 上的有界线性泛函. 该延拓属于 $X^{*}$, 在 $X_{0}$ 上消失而仍然有 $\psi(x_{0})\neq 0$.
\end{proof}

我们将以下推论的证明留作练习.

\begin{corollary}\label{corollary:ch14-稠密判据}
令 $\mathcal{S}$ 为赋范线性空间 $X$ 的子集. 那么 $\mathcal{S}$ 的线性扩张在 $X$ 中稠密当且仅当只要 $\psi\in X^{*}$ 在 $\mathcal{S}$ 上消失则 $\psi=0$.
\end{corollary}

\begin{theorem}\label{theorem:ch14-弱收敛有界}
令 $X$ 为赋范线性空间. 则 $X$ 中的每个弱收敛序列是有界的. 此外, 若在 $X$ 上 $\{x_{n}\}\to x$, 则
\begin{align}
\|x\|\leqslant\liminf\|x_{n}\|.\label{align:ch14-弱收敛范数}
\end{align}
\end{theorem}

\begin{proof}
令在 $X$ 上 $\{x_{n}\}\to x$. 则 $\{J(x_{n}):X^{*}\to\mathbf{R}\}$ 是逐点收敛于 $J(x):X^{*}\to\mathbf{R}$ 的有界泛函序列. 一致有界定理告诉我们 $\{J(x_{n})\}$ 是 $X^{*}$ 上的线性泛函的有界序列. 由于自然嵌入 $J$ 是等距, 序列 $\{x_{n}\}$ 是有界的. 为证明\eqref{align:ch14-弱收敛范数}, 根据定理 7, 存在泛函 $\psi\in X^{*}$ 使得 $\|\psi\|=1$ 而 $\psi(x)=\|x\|$. 现在,
\begin{align*}
\text{对所有 }n,\quad|\psi(x_{n})|\leqslant\|\psi\|\cdot\|x_{n}\|=\|x_{n}\|.
\end{align*}
此外, $\{|\psi(x_{n})|\}$ 收敛到 $|\psi(x)|=\|x\|$. 因此,
\begin{align*}
\|x\|=\lim_{n\to\infty}|\psi(x_{n})|\leqslant\liminf\|x_{n}\|.
\end{align*}
\end{proof}

\textbf{注} Hahn-Banach 定理有相当平易近人的性质. 在它的陈述中需要的数学概念仅仅是线性空间与线性、次可加与正齐次泛函. 除了 Zorn 引理外, 它的证明仅依赖于实数的初等性质. 尽管如此, 通过聪明地选取泛函 $p$, 该定理允许我们为泛函分析创造基本的分析、几何、拓扑的工具. 我们用 Hahn-Banach 定理证明定理 7, 其中 $p$ 选为范数的数乘. 在 14.4 节, 我们用 Hahn-Banach 定理(其中 $p$ 为与某个凸集相伴随的度规泛函)证明线性空间的不交凸子集可被超平面分离. 在下一章, 我们用赋范线性空间到它的二次对偶的自然嵌入证明 Banach 空间 $X$ 的闭单位球是弱序列紧的当且仅当 $X$ 是自反的.

\begin{exercise}\label{exercise:ch14-19}
令 $X$ 为 $C[0,1]$ 的线性子空间, 它是 $L^{2}[0,1]$ 的闭子集. 证明以下断言来证明 $X$ 有有限维数. 序列 $\{f_{n}\}$ 属于 $X$.
\begin{enumerate}[(i)]
\item 证明 $X$ 是 $C[0,1]$ 的闭子空间.
\item 证明存在常数 $M\geqslant 0$, 使得对所有 $f\in X$, 我们有 $\|f\|_{2}\leqslant\|f\|_{\infty}$ 与 $\|f\|_{\infty}\leqslant M\cdot\|f\|_{2}$.
\item 证明对每个 $y\in[0,1]$, 存在 $L^{2}$ 中的函数 $k_{y}$, 使得对每个 $f\in X$, 我们有 $f(y)=\int_{0}^{1}k_{y}(x)f(x)\,dx$.
\item 证明: 若在 $L^{2}$ 弱意义下 $\{f_{n}\}\to f$, 则在 $[0,1]$ 上逐点地 $\{f_{n}\}\to f$.
\item 证明在 $L^{2}$ 弱意义下 $\{f_{n}\}\to f$, 则 $\{f_{n}\}$ 是有界的(在什么意义下?), 因此根据 Lebesgue 控制收敛定理, 在 $L^{2}$ 强意义下 $\{f_{n}\}\to f$.
\item 得出具有范数 $\|\cdot\|_{2}$ 的 $X$ 有紧的闭单位球. 因此, 根据 Riesz 定理, 它是有限维的.
\end{enumerate}
\end{exercise}

\begin{exercise}\label{exercise:ch14-20}
令 $X$ 为赋范线性空间, $\psi$ 属于 $X^{*}$ 而 $\{\psi_{n}\}$ 属于 $X^{*}$. 证明: 若 $\{\psi_{n}\}$ 弱$^{*}$收敛于 $\psi$, 则
\begin{align*}
\|\psi\|\leqslant\liminf\|\psi_{n}\|.
\end{align*}
\end{exercise}

\begin{exercise}\label{exercise:ch14-21}
为 $X=\mathbf{R}^{n}$ 赋予欧氏范数, $Y$ 是 $X$ 的子空间, 而 $\psi:Y\to\mathbf{R}$ 是线性泛函. 定义 $Y^{\perp}$ 为由正交于 $Y$ 的向量组成的 $\mathbf{R}^{n}$ 的线性子空间. 则 $\mathbf{R}^{n}=Y\oplus Y^{\perp}$. 对 $x=y+y'$, $y\in Y$, $y'\in Y^{\perp}$, 定义 $\psi(x)=\psi(y)$. 证明该性质定义了 $\psi\in(\mathbf{R}^{n})^{*}$, 它是 $\psi$ 到 $Y$ 上的延拓, 且与 $\psi|_{Y}$ 有相同的范数.
\end{exercise}

\begin{exercise}\label{exercise:ch14-22}
令 $X=L^{p}=L^{p}[a,b]$, $1<p<\infty$, 而 $m$ 是 Lebesgue 测度. 给定 $L^{p}$ 中的 $f\neq 0$, 对所有 $h\in L^{p}$, 定义
\begin{align*}
\psi(h)=\frac{1}{\|f\|_{p}^{p-1}}\int_{[a,b]}\operatorname{sgn}(f)\cdot|f|^{p-1}\cdot h\,dm.
\end{align*}
用 Hölder 不等式证明 $\psi\in(L^{p})^{*}$, $\|\psi\|=1$ 且 $\psi(f)=\|f\|_{p}$.
\end{exercise}

\begin{exercise}\label{exercise:ch14-23}
对赋范线性空间 $X$ 中的每个点 $x$, 证明
\begin{align*}
\|x\|=\sup\{\psi(x)\mid\psi\in X^{*},\|\psi\|\leqslant 1\}.
\end{align*}
\end{exercise}

\begin{exercise}\label{exercise:ch14-24}
令 $X$ 为赋范线性空间而 $Y$ 是 $X$ 的闭子空间. 证明对每个 $x_{0}\in X\sim Y$, 存在 $\psi\in X^{*}$ 使得:
\begin{align*}
\|\psi\|=1,\quad\text{在 }Y\text{ 上 }\psi=0\text{ 且 }\psi(x_{0})=d,\quad\text{其中 }d=\operatorname{dist}(x_{0},Y)=\inf\{\|x_{0}-y\|\mid y\in Y\}.
\end{align*}
\end{exercise}

\begin{exercise}\label{exercise:ch14-25}
令 $Y$ 为赋范线性空间 $X$ 的线性子空间而 $z$ 是 $X$ 中的向量. 证明
\begin{align*}
\operatorname{dist}(z,Y)=\sup\{\psi(z)\mid\|\psi\|\leqslant 1,\text{在 }Y\text{ 上 }\psi=0\}.
\end{align*}
\end{exercise}

\begin{exercise}\label{exercise:ch14-26}
令 $X$ 为向量空间. $X$ 的子集 $C$ 称为锥, 若只要 $x,y$ 属于 $C$ 而 $\lambda>0$, 则 $x+y\in C$ 且 $\lambda x\in C$. 定义 $X$ 中的偏序 $x\prec y$ 为 $y-x\in C$. $X$ 上的线性泛函 $f$ 称为(关于锥 $C$)是正的, 若在 $C$ 上 $f\geqslant 0$. 令 $Y$ 为 $X$ 的任何具有如下性质的子空间: 对每个 $x\in X$, 存在 $y\in Y$, 使得 $x\prec y$. 假定 $f$ 是 $Y$ 上的线性泛函, 它关于锥 $C\cap Y$ 是正的. 证明 $f$ 可被延拓为 $X$ 上的关于 $C$ 是正的线性泛函. (提示: 采用 Hahn-Banach 引理且用 Zorn 引理找到最大延拓.)
\end{exercise}

\begin{exercise}\label{exercise:ch14-27}
令 $X_{0}$ 为度量空间 $X$ 的子集. 用 Tietze 延拓定理证明, $X_{0}$ 上的每个连续实值函数有到整个 $X$ 的连续延拓当且仅当 $X_{0}$ 是闭的. 这与定理 7 矛盾吗?
\end{exercise}

\begin{exercise}\label{exercise:ch14-28}
令 $(X,\rho)$ 为包含闭集 $F$ 的度量空间. 证明点 $x\in X$ 属于 $F$ 当且仅当 $X$ 上的每个在 $F$ 上消失的连续泛函在 $x$ 也消失. 能用这证明定理 10 吗?
\end{exercise}

\begin{exercise}\label{exercise:ch14-29}
令 $[a,b]$ 为实数的闭有界区间, 考虑 $L^{\infty}[a,b]$, 现在形式上将 $L^{\infty}[a,b]$ 视为关于关系"本性有界函数几乎处处逐点相等"的等价类族. 令 $X$ 为 $L^{\infty}[a,b]$ 的包含连续函数的那些等价类组成的子空间. 证明这样的等价类恰好包含一个连续函数. 因此 $X$ 线性同构于 $C[a,b]$, 模去等同关系, 我们可以考虑 $C[a,b]$ 为 $L^{\infty}[a,b]$ 的线性子空间. 证明 $C[a,b]$ 是 Banach 空间 $L^{\infty}[a,b]$ 的闭子空间.
\end{exercise}

\begin{exercise}\label{exercise:ch14-30}
定义 $\psi:C[a,b]\to\mathbf{R}$ 为对所有 $f\in C[a,b]$, 有 $\psi(f)=f(a)$. 用定理 7 将 $\psi$ 延拓为整个 $L^{\infty}[a,b]$ 上的连续线性泛函(见前一个习题). 证明不存在泛函 $h\in L^{1}[a,b]$ 使得
\begin{align*}
\text{对所有 }f\in L^{\infty}[a,b],\quad\psi(f)=\int_{a}^{b}h\cdot f.
\end{align*}
\end{exercise}

\end{document}